\begingroup\fontsize{10}{14.5}\selectfont\setlength{\tabcolsep}{5pt}\renewcommand{\arraystretch}{1.55}
\textbf{\color{Accent}十二重计数：把 n 个球全部放入 m 个盒子}\par\vspace{5pt}
“可区分”指交换两个球或盒子的身份会成为另一方案。\(S(n,k)\) 数可区分元素的无序非空分组；\(P_m(t)\) 数把整数 \(t\) 无序拆成至多 \(m\) 个正数，重复部件允许。最后一格先给每个非空盒放一个球，再划分剩下的 \(n-m\) 个；方括号为条件成立时取 1，否则取 0。\par\vspace{4pt}
\begin{tabularx}{\linewidth}{ll>{\raggedright\arraybackslash}X>{\raggedright\arraybackslash}X>{\raggedright\arraybackslash}X}\toprule
球 & 盒子 & 可以空着 & 每盒至多一个 & 每盒必须非空\\\midrule
可区分 & 可区分 & \(m^n\) & \(m(m-1)\cdots(m-n+1)\) & \(m!S(n,m)\)\\
可区分 & 不可区分 & \(\sum_{k=0}^{\min(n,m)}S(n,k)\) & \([n\le m]\) & \(S(n,m)\)\\
不可区分 & 可区分 & \(\binom{n+m-1}{m-1}\) & \(\binom mn\) & \(\binom{n-1}{m-1}\)\\
不可区分 & 不可区分 & \(P_m(n)\) & \([n\le m]\) & \(P_m(n-m)\)\\
\bottomrule\end{tabularx}\par\vspace{5pt}
先处理空对象：\(m=0\) 时，只有 \(n=0\) 才有一种方案，十二种口径相同。\(n=0,m>0\) 时，“可以空着”和“每盒至多一个”均为 1，“每盒必须非空”为 0。其余情况下，超范围组合数、负下标划分数均记 0；下降连乘在 \(n>m\) 时记 0，空乘积和 \(0^0\) 都记 1。\par\vspace{5pt}
组合数用阶乘预处理；整行第二类斯特林数用一次卷积；整数划分可按无限件商品处理。令 \(L=\max(n,m)\)，在适用的模数下，这十二个答案一起可 \(O(L\log L)\) 求出。公式的区别来自标签和空盒要求，无需为每一格重新设计 DP。\par\vspace{8pt}
\textbf{\color{Accent}斯特林数：固定元素数看一行，固定组数看一列}\par\vspace{5pt}
\(c(n,k)\) 数 n 元排列恰有 k 个环，称无符号第一类；\(S(n,k)\) 数 n 个可区分元素分成 k 个无序非空组，称第二类。若需要有符号第一类，再乘 \((-1)^{n-k}\)。\par\vspace{4pt}
\begin{tabularx}{\linewidth}{>{\raggedright\arraybackslash}p{28mm}>{\raggedright\arraybackslash}X >{\raggedright\arraybackslash}p{40mm}}\toprule
要数什么 & 为什么这样算 & 批量算法\\\midrule
固定 n，求每种环数 k & 新元素可自成一环，或插到旧环中；答案是连续一次式 \(x(x+1)\cdots(x+n-1)\) 的系数。 & 只求半段再平移倍增，\(O(n\log n)\)。\\
固定 n，求每种非空组数 k & 先给组编号，容斥删去空组，再除组标签数；可拆成 \(j^n/j!\) 与 \((-1)^j/j!\) 两列卷积。 & 一次卷积，含快速幂 \(O(n\log n)\)。\\
固定 k 个环，求所有 \(n\le N\) & 一个大小 t 的环有 \((t-1)!\) 种排列；除去元素标签后，单环数组为 \(a_0=0,a_t=1/t\)。 & 求 k 次幂，第 n 项乘 \(n!/k!\)；Log／Exp，\(O(N\log N)\)。\\
固定 k 个非空组，求所有 \(n\le N\) & 一个大小 t 的集合只有一种内部结构；除去元素标签后，单组数组为 \(a_0=0,a_t=1/t!\)。 & 求 k 次幂，第 n 项乘 \(n!/k!\)；Log／Exp，\(O(N\log N)\)。\\
\bottomrule\end{tabularx}\par\vspace{5pt}
两类都约定第 \((0,0)\) 项为 1；\(k>n\) 或 \(n>0,k=0\) 时为 0。固定列求幂时，先移去零前缀；\(k=0\) 直接返回仅第 0 项为 1。整张 \(n,k\le N\) 表本身有 \(\Theta(N^2)\) 项，用递推 \(c(n,k)=c(n-1,k-1)+(n-1)c(n-1,k)\)、\(S(n,k)=S(n-1,k-1)+kS(n-1,k)\) 已是 \(\Theta(N^2)\)；只算一行／列才有近线性收益。\par\vspace{5pt}
\vspace{5pt}\textbf{\color{Accent}任意定向链：可直接还原的卷积递推}\par
令 \(s_i\) 为第 i 与第 i+1 个位置间的大小号，\(d_i\) 数前 i 个符号中有多少个大于号，\(d_0=0\)。令 \(f_i\) 为前 i 个位置的合法相对排列数除以 \(i!\)，并设 \(f_0=1\)；除阶乘是为了让“给各段分配标签”的组合数从转移中消失，最后再乘回 \(n!\)。\par
按最后一个递增段的起点 j 分类，前面部分由 \(f_j\) 负责，最后一段长 \(i-j\)，贡献 \(1/(i-j)!\)。只能在原有大于号处断段，段内每个被改成小于的大于号给一个负号，因此 \(f_i=\sum_{0\le j<i,\ j=0\ \mathrm{or}\ s_j=\texttt{>}}f_j(-1)^{d_{i-1}-d_j}/(i-j)!\)。\par
转移权的符号可拆成只看目标 i 的 \((-1)^{d_{i-1}}\)，以及只看来源 j 的 \((-1)^{d_j}\)；剩余系数只看距离 \(i-j\)。因此让已求出的 \(f_j(-1)^{d_j}[j=0\text{ 或 }s_j=\texttt{>}]\) 与 \(1/t!\) 卷积，就能批量推进后面的状态；CDQ 保证使用来源时它已算好，总时间 \(O(n\log^2n)\)，答案为 \(n!f_n\)。
\endgroup
